\begin{ouverture}[mouvement]{13}{Inertie et chute libre}{Mouvement et interactions}{Observer · Modéliser · Expérimenter · Expliquer}
\objectif{Relier compensation des forces et vecteur vitesse constant.}
\objectif{Exploiter la contraposée du principe d’inertie.}
\objectif{Analyser une chute verticale avec ou sans vitesse initiale.}
\prerequis{Proportionnalité, conversions et lecture de graphiques.}
\end{ouverture}
\essentiel
\essSection{1}{Le principe d’inertie}
\begin{loi}[Dans un référentiel galiléen]Un système modélisé par un point matériel conserve un vecteur vitesse constant si les forces extérieures se compensent, et réciproquement. Il reste immobile ou en mouvement rectiligne uniforme. Le référentiel terrestre est souvent assimilé à un référentiel galiléen pour les expériences scolaires de courte durée.\end{loi}
Une somme de forces nulle ne signifie pas absence de forces. Un livre immobile subit son poids et l'action du support ; un palet qui glisse sans frottement sur un plan horizontal peut aussi avoir une somme des forces nulle.
\schema[corrige]{inertie}
\essSection{2}{Quand la vitesse change}
Si le vecteur vitesse varie, les forces ne se compensent pas dans les conditions du principe. Sa valeur, sa direction ou les deux peuvent changer. Un mouvement circulaire à valeur de vitesse constante n'est donc pas un mouvement rectiligne uniforme.
\begin{attention}[Une photographie ne suffit pas]Une vitesse nulle à un instant ne prouve pas l'équilibre. Au sommet d'un lancer vertical, la vitesse est momentanément nulle mais le poids reste présent. Pour conclure à l'immobilité, il faut une observation sur une durée.\end{attention}
En Seconde, on relie la variation du vecteur vitesse à une somme non nulle des forces, sans imposer le calcul par une loi d'accélération.
\essSection{3}{Chute libre verticale}
La chute libre est un modèle où seul le poids agit : frottements de l'air et poussée sont négligés. Elle peut commencer avec une vitesse initiale vers le haut, vers le bas, ou nulle.
\schema[corrige]{chute}
Lors d'une descente, le vecteur vitesse est vers le bas et sa valeur augmente. Lors d'une montée verticale sous le seul poids, le vecteur vitesse est vers le haut mais sa valeur diminue ; au sommet, il devient nul puis s'oriente vers le bas. Le poids reste orienté vers le bas durant les trois phases.
\essSection{4}{Expérimenter et modéliser}
\begin{experience}[Tester le modèle de chute]Filmer une petite balle devant une échelle verticale ; relever des positions à temps égaux et représenter les vitesses approchées. Comparer avec une feuille lâchée à plat : si les effets de l'air sont importants, le modèle de chute libre n'est plus adapté. Délimiter la zone et utiliser une hauteur sûre, sous encadrement.\end{experience}
\exercices
\rubrique{Prendre confiance}
\begin{blocLivre}{ex-13-1}\begin{exercice}[1]{Un livre immobile}
Un livre est posé sur une table horizontale.
\schema{inertie}
\begin{questions}
\item Les forces sur le livre se compensent-elles ?
\item Cela signifie-t-il qu’aucune force n’agit ?
\end{questions}
\end{exercice}\end{blocLivre}
\begin{blocLivre}{ex-13-2}\begin{exercice}[1]{Un palet idéal}
Un palet glisse horizontalement ; les frottements sont négligés et les forces verticales se compensent.
\begin{questions}
\item Prévoir son mouvement.
\item Faut-il une force horizontale pour entretenir cette vitesse ?
\end{questions}
\end{exercice}\end{blocLivre}
\begin{blocLivre}{ex-13-3}\begin{exercice}[1]{Une voiture qui freine}
La valeur de la vitesse d’une voiture diminue sur une droite.
\begin{questions}
\item Le vecteur vitesse est-il constant ?
\item Les forces se compensent-elles ?
\end{questions}
\end{exercice}\end{blocLivre}
\begin{blocLivre}{ex-13-4}\begin{exercice}[1]{Définir la chute libre}
Une balle descend verticalement et les effets de l’air sont négligés.
\begin{questions}
\item Quelle force reste dans le modèle ?
\item La vitesse doit-elle être nulle au départ ?
\end{questions}
\end{exercice}\end{blocLivre}
\rubrique{S’entraîner}
\begin{blocLivre}{ex-13-5}\begin{exercice}[2]{Forces compensées}
Un objet subit deux forces de 3,0 N, verticales et opposées, et aucune autre force.
\begin{questions}
\item Calculer leur somme vectorielle.
\item Peut-on choisir entre repos et mouvement sans autre information ?
\end{questions}
\end{exercice}\end{blocLivre}
\begin{blocLivre}{ex-13-6}\begin{exercice}[2]{Au sommet du lancer}
Une balle est lancée verticalement vers le haut. On néglige l’air.
\schema{chute}
\begin{questions}
\item Comparer vitesse et poids au sommet.
\item Pourquoi ce point n’est-il pas un état d’équilibre ?
\end{questions}
\end{exercice}\end{blocLivre}
\begin{blocLivre}{ex-13-7}\begin{exercice}[2]{Un cercle uniforme}
Un objet décrit un cercle à valeur de vitesse constante.
\schema{mouvement-circulaire}
\begin{questions}
\item Identifier ce qui varie dans le vecteur vitesse.
\item Que conclure sur la somme des forces ?
\end{questions}
\end{exercice}\end{blocLivre}
\begin{blocLivre}{ex-13-8}\begin{exercice}[2]{Une vitesse limite}
Une bille descend dans un liquide, en ligne droite à vitesse constante après une phase transitoire.
\begin{questions}
\item Les forces se compensent-elles pendant cette phase uniforme ?
\item Peut-on la modéliser par la seule action du poids ?
\end{questions}
\end{exercice}\end{blocLivre}
\rubrique{Raisonner et approfondir}
\begin{blocLivre}{ex-13-9}\begin{exercice}[3]{Interpréter des vecteurs}
À trois instants d’une descente verticale, les vitesses valent 1,0 ; 2,0 ; 3,0 m/s vers le bas.
\begin{questions}
\item Décrire la variation du vecteur vitesse.
\item Cette observation est-elle compatible avec une somme des forces nulle ?
\end{questions}
\end{exercice}\end{blocLivre}
\begin{blocLivre}{ex-13-10}\begin{exercice}[3]{Une balle montante}
Une balle monte verticalement après avoir quitté la main ; l’air est négligé.
\begin{questions}
\item Existe-t-il encore une force de la main sur la balle ?
\item Pourquoi la balle ralentit-elle sans que son poids change de sens ?
\end{questions}
\end{exercice}\end{blocLivre}
\begin{blocLivre}{ex-13-11}\begin{exercice}[3]{Choisir le système}
Une personne debout sur un sol horizontal est immobile.
\begin{questions}
\item Le poids et l’action du sol se compensent-ils sur la personne ?
\item Sont-ils une paire d’actions réciproques ?
\end{questions}
\end{exercice}\end{blocLivre}
\begin{blocLivre}{ex-13-12}\begin{exercice}[3]{Un diagnostic trop rapide}
À un instant, un mobile se déplace vers la droite et la somme de ses forces est non nulle.
\begin{questions}
\item Peut-on en déduire qu’il va nécessairement plus vite ?
\item Quelles évolutions du vecteur vitesse sont possibles selon les actions ?
\end{questions}
\end{exercice}\end{blocLivre}

